%%%%%%%%%%%%%%%%%%%%%%%%%%%%%%%%%%%%%%%%%%%%%%%%%%%%%%%%%%
%   Autoren:
%   Prof. Dr. Bernhard Drabant
%   Prof. Dr. Dennis Pfisterer
%   Prof. Dr. Julian Reichwald
%%%%%%%%%%%%%%%%%%%%%%%%%%%%%%%%%%%%%%%%%%%%%%%%%%%%%%%%%%

%%%%%%%%%%%%%%%%%%%%%%%%%%%%%%%%%%%%%%%%%%%%%%%%%%%%%%%%%%
%	ANLEITUNG: 
%   1. Ersetzen Sie firmenlogo.jpg im Verzeichnis img
%   2. Passen Sie alle Stellen im Dokument an, die mit 
%      @stud 
%      markiert sind
%%%%%%%%%%%%%%%%%%%%%%%%%%%%%%%%%%%%%%%%%%%%%%%%%%%%%%%%%%

%%%%%%%%%%%%%%%%%%%%%%%%%%%%%%%%%%%%%%%%%%%%%%%%%%%%%%%%%%
%	ACHTUNG: 
%   Für das Erstellen des Literaturverzeichnisses wird das 
%   modernere Paket biblatex in Kombination mit biber 
%   verwendet - nicht mehr das ältere Paket BibTex!
%
%   Bitte stellen Sie Ihre TeX-Umgebung entsprechend ein (z.B. TeXStudio): 
%   Einstellungen --> Erzeugen --> Standard Bibliographieprogramm: biber
%%%%%%%%%%%%%%%%%%%%%%%%%%%%%%%%%%%%%%%%%%%%%%%%%%%%%%%%%%

\documentclass[fontsize=12pt,BCOR=5mm,DIV=12,parskip=half,listof=totoc,
               paper=a4,toc=bibliography,pointlessnumbers]{scrreprt}
               
               %toc=listof,listof=entryprefix,
               
\makeindex

%% Elementare Pakete, Konfigurationen und Definitionen werden geladen (gegebenenfalls anpassen)
% !TEX root =  master.tex

%%%%%%%%%%%%%%%%%%%%%%%%%%%%%%%%%%%%%%%%%%%%%%%%%%%%%%%%%%%%%%%%%%
%	ANLEITUNG: 
% Passen Sie gegebenenfalls alle Stellen im Dokument an, die mit 
% @stud 
% markiert sind.
%%%%%%%%%%%%%%%%%%%%%%%%%%%%%%%%%%%%%%%%%%%%%%%%%%%%%%%%%%%%%%%%%%

\usepackage{makeidx}                  % allows index generation
\usepackage{listings}	                %Format Listings properly
\usepackage{lipsum}                   % Blindtext
\usepackage{graphicx}                 % use various graphics formats
\usepackage[german]{varioref}         % nicer references \vref
\usepackage{caption}	                % better Captions
\usepackage{booktabs}                 % nicer Tabs
\usepackage[hidelinks=true]{hyperref} % keine roten Markierungen bei Links
\usepackage{fnpct}                    % Correct superscripts 
\usepackage{calc}                     % Used for extra space below footsepline, in particular
\usepackage{array}
\usepackage{acronym}
\usepackage{algorithm}
\usepackage{algpseudocode}
\usepackage{setspace}
\usepackage{tocloft}
\usepackage{float}

%% Schriftarten- und Zeichenpakete
\usepackage[T1]{fontenc}
\usepackage[utf8]{inputenc}

%%
%% @stud
%%
%%	FONT SELECTION: Schriftarten und Schriftfamilie
%%%%%%%%%%%%%
%% SCHRIFTART
%%%%%%%%%%%%%
% 0) without decomment: normal font families 
% ...
% 1) Latin Modern 
%\usepackage{lmodern}        
% 2) Times 
%\usepackage{mathptmx}         
% 3) Helvetica
%\usepackage[scaled=.92]{helvet} 
%%%%%%%%%%%%%%%%%%
%%	SCHRIFTFAMILIE
%%%%%%%%%%%%%%%%%%
% ohne Serifen
\renewcommand*{\familydefault}{\sfdefault}
\addtokomafont{disposition}{\sffamily}
%
% mit Serifen
%\renewcommand*{\familydefault}{\rmdefault}
%\addtokomafont{disposition}{\rmfamily}
%
% Typewriter
%\renewcommand*{\familydefault}{\ttdefault}
%\addtokomafont{disposition}{\ttfamily}

%%
%% @stud
%%
%% LANGUAGE SETTINGS
\usepackage[ngerman]{babel} 	        % german language
\usepackage[german=quotes]{csquotes} 	% correct quoting using \enquote{}
%\usepackage[english]{babel}          % english language
%\usepackage{csquotes} 	              % correct quoting using \enquote{}

%%
%% @stud
%%
%% Uncomment the following lines to support hard URL breaks in bibliography 
%\apptocmd{\UrlBreaks}{\do\f\do\m}{}{}
%\setcounter{biburllcpenalty}{9000}% Kleinbuchstaben
%\setcounter{biburlucpenalty}{9000}% Großbuchstaben

%%
%% @stud
%%
%% FOOTNOTES: Count footnotes over chapters
%% \counterwithout{footnote}{chapter}

%	ACRONYMS
\makeatletter
\@ifpackagelater{acronym}{2015/03/20}
{\renewcommand*{\aclabelfont}[1]{\textbf{{\acsfont{#1}}}}}{}
\makeatother

%	LISTINGS
% @stud: ggf. Namen/Text anpassen (englisch)
\renewcommand{\lstlistingname}{Quelltext} 
\renewcommand{\lstlistlistingname}{Quelltextverzeichnis}
\lstset{numbers=left,
	numberstyle=\tiny,
	captionpos=b,
	basicstyle=\ttfamily\small}

%	ALGORITHMS
% @stud: ggf. Namen/Text anpassen (englisch)
\renewcommand{\listalgorithmname}{Algorithmenverzeichnis}
\floatname{algorithm}{Algorithmus}

%	PAGE HEADER / FOOTER
%	Warning: There are some redefinitions throughout the master.tex-file!  DON'T CHANGE THESE REDEFINITIONS!
\RequirePackage[automark]{scrlayer-scrpage}
%alternatively with separation lines: \RequirePackage[automark,headsepline,footsepline]{scrlayer-scrpage}

\renewcommand{\chaptermarkformat}{}
\RedeclareSectionCommand[beforeskip=0pt]{chapter}
\clearscrheadfoot

%\ifoot[\rule{0pt}{\ht\strutbox+\dp\strutbox}DHBW Mannheim]{\rule{0pt}{\ht\strutbox+\dp\strutbox}DHBW Mannheim}
\ofoot[\rule{0pt}{\ht\strutbox+\dp\strutbox}\pagemark]{\rule{0pt}{\ht\strutbox+\dp\strutbox}\pagemark}
\ohead{\headmark}

\newcommand{\TitelDerArbeit}[1]{\def\DerTitelDerArbeit{#1}\hypersetup{pdftitle={#1}}}
\newcommand{\AutorDerArbeit}[1]{\def\DerAutorDerArbeit{#1}\hypersetup{pdfauthor={#1}}}
\newcommand{\Firma}[1]{\def\DerNameDerFirma{#1}}
\newcommand{\Kurs}[1]{\def\DieKursbezeichnung{#1}}
\newcommand{\Abteilung}[1]{\def\DerNameDerAbteilung{#1}}
\newcommand{\Studiengangsleiter}[1]{\def\DerStudiengangsleiter{#1}}
\newcommand{\WissBetreuer}[1]{\def\DerWissBetreuer{#1}}
\newcommand{\FirmenBetreuer}[1]{\def\DerFirmenBetreuer{#1}}
\newcommand{\Bearbeitungszeitraum}[1]{\def\DerBearbeitungszeitraum{#1}}
\newcommand{\Abgabedatum}[1]{\def\DasAbgabedatum{#1}}
\newcommand{\Matrikelnummer}[1]{\def\DieMatrikelnummer{#1}}
\newcommand{\Studienrichtung}[1]{\def\DieStudienrichtung{#1}}
\newcommand{\ArtDerArbeit}[1]{\def\DieArtDerArbeit{#1}}
\newcommand{\Literaturverzeichnis}{Literaturverzeichnis}

\newcommand{\settingBibFootnoteCite}{
	\setlength{\bibparsep}{\parskip}		  % Add some space between biblatex entries in the bibliography
	\addbibresource{bibliography.bib}	    % Add file bibliography.bib as biblatex resource
	\DefineBibliographyStrings{ngerman}{andothers = {{et\,al\adddot}},}
}

\newcommand{\setTitlepage}{
	% !TEX root =  master.tex
% @stud: ggf. Namen/Text anpassen (englisch)
\begin{titlepage}
\begin{minipage}{\textwidth}
		\vspace{-2cm}
		\noindent \includegraphics[scale=0.09]{\imagedir/firmenlogo.jpg} \hfill \includegraphics{\imagedir/logo.jpg}
\end{minipage}
\vspace{1em}
%\sffamily
\begin{center}
	{\textsf{\large Duale Hochschule Baden-W\"urttemberg Mannheim}}\\[4em]
	{\textsf{\textbf{\large{\DieArtDerArbeit}report}}}\\[6mm]
	{\textsf{\textbf{\Large{}\DerTitelDerArbeit}}} \\[1.5cm]
	{\textsf{\textbf{\large{}Studiengang Wirtschaftsinformatik}}\\[6mm]
	\textsf{\textbf{Studienrichtung \DieStudienrichtung}}}\vspace{10em}
	
\\[3.5cm]
	
	\begin{minipage}{\textwidth}
		\begin{tabbing}
		Wissenschaftliche(r) Betreuer(in): \hspace{0.85cm}\=\kill
		Verfasser: \> \DerAutorDerArbeit \\[1.5mm]
		Matrikelnummer: \> \DieMatrikelnummer \\[1.5mm]
		Firma: \> \DerNameDerFirma  \\[1.5mm]
		Modul: \> \DerNameDerAbteilung \\[1.5mm]
		Kurs: \> \DieKursbezeichnung \\[1.5mm]
		Dozent: \> \DerStudiengangsleiter \\[1.5mm]
		
		
		Bearbeitungszeitraum: \> \DerBearbeitungszeitraum\\[1.5mm]
%		alternativ:\\[1.5mm]
%		Eingereicht: \> \DasAbgabedatum	
		\end{tabbing}
	\end{minipage}
\end{center}
\end{titlepage}
	\pagenumbering{roman} % Römische Seitennummerierung
	\normalfont	
}

\newcommand{\initializeText}{
	\clearpage
	\ihead{\chaptername~\thechapter} % Neue Header-Definition
	\pagenumbering{arabic}           % Arabische Seitenzahlen
}

\newcommand{\initializeBibliography}{
	\ihead{}
	\printbibliography[title=\Literaturverzeichnis] 
	\cleardoublepage
}

\newcommand{\initializeAppendix}{
	\appendix
  \ihead{}
  \cftaddtitleline{toc}{chapter}{Anhang}{}
}



%%
%% @stud
%%
%% PERSÖNLICHE ANGABEN (BITTE VOLLSTÄNDIG EINGEBEN zwischen den Klammern: {...})
%%
\ArtDerArbeit{Projekt} % "Bachelor" oder "Projekt" wählen
\TitelDerArbeit{Bundesliga Match Predictions}
\AutorDerArbeit{Max Bernauer, Philipp Dingfelder, Julius Könning}
\Abteilung{Data Exploration Project}
\Firma{SAP SE, Schaeffler}
\Kurs{WWI20DSB}
\Studienrichtung{Data Science}
\Matrikelnummer{5763624, 8687786, 7305370}
\Studiengangsleiter{Simon Poll}
\WissBetreuer{x}
\FirmenBetreuer{x}
\Bearbeitungszeitraum{Sommersemester 2022}
\Abgabedatum{31.07.2022}

%%
%% @stud
%%
%% BIBLIOGRAPHY (@stud: Bibliographie-Stil wählen - Position und Indizierung)
%%  Auswahl zwischen: 
%%   NUMERIC Style
%%   IEEE Style
%%   ALPHABETIC Style
%%   HARVARD Style
%%   CHICAGO Style 
%%   (oder eigenen zulässigen Stil wählen) 
%%
%%%%%%%%%%%%%
%% Zitierstil
%%%%%%%%%%%%%
% NUMERIC Style - e. g. [12]
%\newcommand{\indextype}{numeric} 
%
% IEEE Style - numeric kind of style 
%\newcommand{\indextype}{ieee} 
%
% ALPHABETIC Style - e. g. [AB12]
%\newcommand{\indextype}{alphabetic} 
%
% HARVARD Style 
\newcommand{\indextype}{apa} 
%
% CHICAGO Style 
%\newcommand{\indextype}{authoryear}
%
%%%%%%%%%%%%%%%%%%%%%%
%% Position des Zitats
%%%%%%%%%%%%%%%%%%%%%%
\newcommand{\position}{inline} 
%
% (!!) FOOTNOTE POSITION NOT RECOMMENDED IN MINT DOMAIN:
%\newcommand{\position}{footnote}

%% Final: Setzen des Zitierstils und der Zitatposition

\usepackage[backend=biber, autocite=\position, style=\indextype]{biblatex} 	
%\settingBibFootnoteCite

% Citavi einbinden
\bibliography{bibliography}

%%
%% Definitionen und Commands
%%
\newcommand{\abs}{\par\vskip 0.2cm\goodbreak\noindent}
\newcommand{\nl}{\par\noindent}
\newcommand{\mcl}[1]{\mathcal{#1}}
\newcommand{\nowrite}[1]{}
\newcommand{\NN}{{\mathbb N}}
\newcommand{\imagedir}{img}

\makeindex

\begin{document}

\setTitlepage

%%%%%%%%%%%%%%%%%%%%%%%%%%%%%%%%%%%
% EHRENWÖRTLICHE ERKLÄRUNG
%
% @stud: ewerkl.tex bearbeiten
%
% !TEX root =  master.tex
\clearpage
\chapter*{Ehrenwörtliche Erklärung}

% Wird die folgende Zeile auskommentiert, erscheint die ehrenwörtliche
% Erklärung im Inhaltsverzeichnis.

% \addcontentsline{toc}{chapter}{Ehrenwörtliche Erklärung}
Ich versichere hiermit, dass ich die vorliegende Arbeit mit dem Titel ``\textit{\DerTitelDerArbeit}'' selbstständig verfasst und 
keine anderen als die angegebenen Quellen und Hilfsmittel benutzt habe. Ich versichere zudem, dass die eingereichte elektronische 
Fassung mit der gedruckten Fassung übereinstimmt.

\vspace{3cm}
Ort, Datum \hfill \DerAutorDerArbeit
 
\cleardoublepage  
%%%%%%%%%%%%%%%%%%%%%%%%%%%%%%%%%%%

%%%%%%%%%%%%%%%%%%%%%%%%%%%%%%%%%%%
% SPERRVERMERK
%
% @stud: nondisclosurenotice.tex bearbeiten
%
%% !TEX root =  master.tex
\chapter*{Sperrvermerk}

Der Inhalt dieser Arbeit darf weder als Ganzes noch in Auszügen Personen außerhalb des Prüfungsprozesses 
und des Evaluationsverfahrens zugänglich gemacht werden, sofern keine anders lautende Genehmigung der Ausbildungsstätte vorliegt. 

\cleardoublepage
 
%\cleardoublepage
%%%%%%%%%%%%%%%%%%%%%%%%%%%%%%%%%%%

%%%%%%%%%%%%%%%%%%%%%%%%%%%%%%%%%%%
%	KURZFASSUNG
%
% @stud: acknowledge.tex bearbeiten
%
%% !TEX root =  master.tex
%\chapter*{Danksagung}

%Hier können Sie eine Danksagung schreiben. 



%\cleardoublepage 
%%%%%%%%%%%%%%%%%%%%%%%%%%%%%%%%%%%

%%%%%%%%%%%%%%%%%%%%%%%%%%%%%%%%%%%
% VERZEICHNISSE und ABSTRACT
%
% @stud: ggf. nicht benötigte Verzeichnisse auskommentieren/löschen
%
\tableofcontents
\cleardoublepage

% Abbildungsverzeichnis
%\phantomsection
%\addcontentsline{toc}{chapter}{\listfigurename}
%\listoffigures
%\cleardoublepage

%	Tabellenverzeichnis
%\phantomsection
%\addcontentsline{toc}{chapter}{\listtablename}
%\listoftables
%\cleardoublepage

% Algorithmenverzeichnis
%\listofalgorithms
%\cleardoublepage

% Abkürzungsverzeichnis
% @stud: acronyms.tex bearbeiten
%% !TEX root =  master.tex
\clearpage
\chapter*{Abkürzungsverzeichnis}	
\addcontentsline{toc}{chapter}{Abkürzungsverzeichnis}

\begin{acronym}[XXXXXXX]
	\acro{xxx}[XXX]Description}
\end{acronym}

 
%\cleardoublepage

%	Kurzfassung / Abstract
% @stud: abstract.tex bearbeiten
%% !TEX root =  master.tex
\chapter*{Kurzfassung (Abstract)}
\addcontentsline{toc}{chapter}{Kurzfassung (Abstract)}
 
\cleardoublepage

%%%%%%%%%%%%%%%%%%%%%%%%%%%%%%%%%%%%%%%%%%%%%%%%%%%%%%%%%%%%%%%%%%%%%%%%%%%%%%%%%%%%%%%%%%
% KAPITEL UND ANHÄNGE
%
% @stud:
%   - nicht benötigte: auskommentieren/löschen
%   - neue: bei Bedarf hinzufügen mittels input-Kommando an entsprechender Stelle einfügen
%%%%%%%%%%%%%%%%%%%%%%%%%%%%%%%%%%%%%%%%%%%%%%%%%%%%%%%%%%%%%%%%%%%%%%%%%%%%%%%%%%%%%%%%%%

\initializeText
\onehalfspacing

%%%%%%%%%%%%%%%%%%%%%%%%%%%%%%%%%%%
% KAPITEL
%
% @stud: einzelne Kapitel bearbeiten und eigene Kapitel hier einfügen
%
% Einleitung
% !TEX root =  master.tex
\chapter{Zufall und Vorhersehbarkeit im Fußball}
	
Was macht Fußball so spannend und beliebt? Sepp Herberger hat dazu eine klare Meinung:
\glqq Fußball ist deshalb spannend, weil niemand weiß, wie das Spiel ausgeht.\grqq 

Dass Sepp Herberger teilweise Recht hat, bestätigen immer wieder vermeintliche Underdogs, die den klaren Favoriten in einem Spiel schlagen oder sogar düpieren können. Hierfür gibt es Statistiken, die den Effekt erklären: So fallen 47\% aller Tore zufällig, also durch einen Fehler in der Verteidigung oder einen abgefälschten Schuss \autocite[Vgl.][]{.06102018}. Außerdem ist jeder dritte Ballkontakt ein unvorhersehbares Ereignis. Es wird also beinahe unmöglich einen perfekten Spielzug mit anschließendem Torerfolg zu planen \autocite[Vgl.][]{.}.

Dennoch gibt es auch Komponenten, die das Gegenteil vermuten lassen.
So wissen wir aus Berechnungen, dass die Mannschaft mit dem höheren Marktwert zu 48\% gewinnt, das Heimteam zu 45\% oder das Team mit der besseren Tabellenplatzierung zu 44\%. 
Inwieweit ist Fußball also etwa durch Machine Learning vorhersehbar oder überwiegt doch die Zufallskomponente? Und ist es tatsächlich möglich Anbieter wie Bet365 und Bwin durch ein mathematisches Modell zu schlagen? Diese beiden Fragen versuchen wir im Rahmen der Arbeit zu beantworten.

Das erfolgreiche Vorhersagen von Fußballspielen wäre auch mit einem enormen wirtschaftlichen Nutzen versehen. So kann man durch eine solche Prediction potenziell starken Profit bei etwaigen Wettanbietern erzielen. Andererseits kann man diese Software auch an Wettanbieter verkaufen, damit diese ihre Quoten noch effizienter und genauer berechnen können.




% mehrere Grundlagen- und Forschungs-Kapitel
% !TEX root =  master.tex
\chapter{Theoretische Grundlagen}

\section{Related Work}
Nach Recherche von Vergleichsprojekten wurden einige Ideen aufgegriffen. Beim Preprocessing wurden einige relevante Features hinzugefügt. Begonnen wurde hierbei mit Elo und Angriffs- und Verteidigungsstärke \autocite[Vgl.][]{Accso.14102021}. Als Vergleichsprojekt wurde ein mathematischer Ansatz gefunden. Dieser ist mit R programmiert und schafft laut angaben der Quelle 64\% accuracy \autocite[Vgl.][]{Doan.15.3.2019}. Für die grundlegende Idee des Business Use Case wurde eine weitere Quelle herangezogen \autocite[Vgl.][]{Hartley.05102022}.

\section{Verwendete Technologien und Bibliotheken}
Für die Datenvorbereitung werden klassische Python-Module wie Pandas, Sklearn und Matplotlib genutzt.

Bei der Modellimplementierung wurde vorerst ein Decision Tree Classifier genutzt. Dieser stammt aus dem Sklearn-Modul tree und erreichte keine nennenswerten Ergebnisse. Nach Ausprobieren dieses Classifiers wurde sich für das Ensemble-Learning entschieden. Dabei nutzt man mehrere Algorithmen um sie zeitgleich zu vergleichen und potenziell zu vereinen. Die Algorithmen wie beispielsweise Logistic Regression oder Gaussian Naive Bayes stammen auch von Sklearn. Final wurde ein Neural Network implementiert. Dazu wurde Keras genutzt. Aus Keras wurde außerdem eine Random Parameter Search zur weiteren Optimierung angewandt. Durch Matplotlib wurden die Ergebnisse visualisiert. Außerdem wurde eine Poisson-Verteilung genutzt um einen mathematischen Ansatz zu implementieren.
% !TEX root =  master.tex
\chapter{Praktischer Teil}

\section{Umsetzung}
- datenbeschaffung und vorverarbeitung
- implementation

\section{Ergebnisse}
- welches modell erzielt welches ergebnis

% Fazit und Ausblick
% !TEX root =  master.tex
\chapter{Fazit}

Abschließend lässt sich festellen, dass sich die Accuracy des Neural Networks
im Verlauf des Projekts Schritt für Schritt durch das optimieren der Datenbasis
und des Modells selber erhöhen ließ. Die größte Erhöhung dabei kam durch das fokussieren
auf den direkten Unterschied der Werte der Mannschaften sowie die Random Paramter Search
des Modells. Final konnte eine 53\% Accuracy erreicht werden.

Wie in der Grafik in Abschnitt drei von Kapitel 3 gezeigt war es möglich mit Hilfe
des Modells profitable zu wetten. Vergleicht man die Profit Kurve es Modells jedoch
mit der der Wettanbieter fällt auf, dass diese im Verlauf der Saison Intervalle mit
höheren Gewinnen aufweisen, ihre Predictions also besser funktionieren. Das lässt sich
darauf zurück führen, dass das Modell mit maximal 65\% Sicherheit das Ergebnis eines Spiels
vorhersagt. Die Modelle der Wettanbieter weißen eine durchschnittlich höhere Sicherheit in Bezug auf die
Ergebnisse der Spiele aus, was dazu führt, dass diese in der Simulation höhere Beträge
einsetzten können und insgesamt auch mehr wetten eingehen als unser Modell.
Um also langfristig die Wettanbieter gewinnbringend zu schlagen müsste das Modell
weiterverbessert werden um eine höhere Accuracy zu erhalten. Mögliche Ansätze hierfür
sind:

Zusammenfassend lassen sich die Ergebnisse des Projekts folgendermaßen zusamenfassen:
Das Projekt war erfolgreich. Durch die 53 prozentige Accuracy sowie das halbwegs
profitbale Wetten konnte gezeigt werden, dass sich Fußballergebnisse zumindest in
Ansätzen algorithmisch mit Hilfe von Modellen beschreiben lassen. Damit ist der Use
Case zum Großteil erfüllt. Es ist jedoch kritisch anzumerken, dass auch sehr viele
nicht genau spezifizierbare Zufallsvariablen einen Einfluss auf den Ausgang
eines Spiels haben, sodass sich nie mit 100\%er Wahrscheinlichkeit der genaue Ausgang
eines Fußballspiels vorhersagen lässt.


Die Modell Accuracy von 53\% ist für sich genommen relativ vielversprechend.
Sie zeigt, dass sich Fußballergebnisse zumindest in Ansätzen algorithmisch
bestimmen lassen. Hierbei ist jedoch kritisch anzumerken, dass auch sehr viele nicht genau
spezifizierbare Zufallsvariablen einen Einfluss auf den Ausgang eines Spiels haben.



Abschließend lässt sich feststellen, dass man durch die
Optimierung des Neural Networks die höchste Accuracy erreicht.
Nach der Random Parameter Search konnten wir die Accuracy erhöhen.
Wie bereits erwähnt lag diese final bei 53 Prozent.


Mithilfe dieses Modells wurde das Betting bei Wettanbietern simuliert. Leider
 kann man sehen, dass unser Modell schnell negative Summen erreicht
  und einiges an Geld verliert. Der Fußball ist wohl doch ein
  Gebiet mit zu vielen Zufallsvariablen, sodass man keine
  Genauigkeit erzielen kann, die einen Wettanbieter schlägt.


  Die Genauigkeit von 53 Prozent ist jedoch für sich genommen
  recht vielversprechend. Es zeigt, dass sich der Fußball in
  Ansätzen algorithmisch bestimmen lässt. Für unseren Business
  Use Case ist dieses Ergebnis natürlich nicht hoch genug und
  somit muss dieser verworfen werden. Es hat sich gezeigt, dass
  man schnell Geld verliert bei der Nutzung unseres Modells.

%%%%%%%%%%%%%%%%%%%%%%%%%%%%%%%%%%%

%%%%%%%%%%%%%%%%%%%%%%%%%%%%%%%%%%%
% ANHÄNGE
%
% @stud: einzelne Anhänge bearbeiten und eigene Anhänge hier einfügen 
%        die nachfolgenden Zeilen deaktivieren, wenn keine Anhänge verwendet werden
%
\initializeAppendix
% !TEX root =  master.tex
\chapter{Beschreibung des Quellcodes}
\section{Ausführen des Programms}
\begin{enumerate}
	\item Klonen des Github-Repository unter dem Link https://github.com/Max280201/Data-Exploration-Project (Branch: main)
	\item Öffnen der CMD und navigieren in den Git-Hub-Ordner
	\item CMD-Command 'py -3.X -m venv Hoyzer' (Bash: 'python3.X -m venv HoyzerVenv') ausführen (erstellt die virtuell Environment, ersetzen von X durch die vorhandene Python-Version)
	\item \grqq HoyzerVenv/Scripts/activate.bat\grqq in CMD ausführen (aktiviert die virtuell Environment)
	\item CMD-Command \grqq pip install -r requirements.txt\grqq{} ausführen
	\item Auswählen der python version der venv als Kernel; Ausführen der Python-Datei/ Notebooks anschließend möglich
\end{enumerate}

\section{Erklärung der einzelnen Bestandteile}
\begin{itemize}
	\item preprocessing\_pipeline\_v1.ipynb: kombiniert die csv-Dateien der Saisonspiele und die Marktwerte; außerdem Datenvorverarbeitung und Berechnung weiterer Features
	\item modelltraining\_prod\_v4.py: Modelltraining eines NN (sowohl eines, um die letzte Saison vohrerzusagen, als auch eines, um zufällige Spiele des Bereichs vorherzusagen) und kurze Evaluation des Modells
	\item modelltraining\_test\_area\_v4.ipynb: enthält alle relevanten getesteten Techniken für das Modell und Alternativen (Hyperparametertuning, Ensemble-Learning, Poission-Prediction);
	
	Anmerkung: Random Search kann zu dem Fehler \grqq Access denied\grqq{} führen, dies lässt sich durch eine geringere Anzahl an Wiederholungen beheben (https://github.com/keras-team/keras-tuner/issues/339)
	
	\item translate\_betting\_odds.ipynb: Auswertung und Vorverarbeitung der Wettquoten
	\item evaluate\_model\_odd\_predictions\_v2.ipynb: Vergleich des trainierten Modells mit Wettanbietern und Auswertung
\end{itemize}
%%%%%%%%%%%%%%%%%%%%%%%%%%%%%%%%%%%

\singlespacing

%%%%%%%%%%%%%%%%%%%%%%%%%%%%%%%%%%%
% LITERATURVERZEICHNIS
% @stud: Literaturverzeichnis in Datei bibliography.bib anpassen. 
%
% Alternative zu Verwendung von \initializeBibliography: Citavi ...
% (dann \initializeBibliography auskommentieren und eigenes LaTex Coding verwenden)
%

\printbibliography
%%%%%%%%%%%%%%%%%%%%%%%%%%%%%%%%%%%

%%%%%%%%%%%%%%%%%%%%%%%%%%%%%%%%%%%
% INDEX
% @stud: ggf. Index auskommentieren, wenn nicht benötigt
%
%\addcontentsline{toc}{chapter}{Index}
%\printindex

\end{document}
